% Part: model-theory
% Chapter: models-of-arithmetic
% Section: computable-models

\documentclass[../../../include/open-logic-section]{subfiles}

\begin{document}

\olfileid{mod}{mar}{cmp}
\section{अंकगणिताची संगणनक्षम प्रतिमाने}

\begin{explain}
मानक प्रतिमान~$\Struct{N}$ ला दोन उपयुक्त वैशिष्ट्ये आहेत. त्याचा प्रांत
नैसर्गिक संख्यांचा संच~$\Nat$ आहे, म्हणजे त्याचे घटक हे नेमके तेच प्रकारचे
विषय आहेत ज्यांच्याविषयी अंकगणिताची भाषा वापरून आपल्याला बोलायचे असते, आणि
मानक संख्यांक~$\num{n}$ हा खरोखर~$n$ लाच निर्देशित करतो. दुसरे उपयुक्त
वैशिष्ट्य म्हणजे~$\Lang{L_A}$ मधील अतार्किक चिन्हांची सर्व अर्थनिर्धारणे
\emph{संगणनक्षम} आहेत. $\Assign{\prime}{N}$, $\Assign{+}{N}$ आणि
$\Assign{\times}{N}$ म्हणून काम करणारी उत्तरवर्ती, बेरीज आणि गुणाकार ही
फलने संख्यांवरील संगणनक्षम फलने आहेत. (म्हणजे ट्यूरिंग यंत्रांनी संगणन करता
येणारी किंवा, उदाहरणार्थ, आदिम पुनरावर्तनाने परिभाषित करता येणारी.) तसेच
$\Struct{N}$ वरील लघुतर संबंध, म्हणजे~$\Assign{<}{N}$, निर्णेय आहे.

$\Th{Q}$ आणि $\Th{PA}$ यांसारख्या अंकगणितीय उपपत्तींच्या अमानक प्रतिमानांत
अमानक घटक असलेच पाहिजेत. त्यामुळे त्यांच्या प्रांतांत सामान्यतः~$\Nat$
व्यतिरिक्तही !!{element}s असतात. तरी कोणतीही गणनीय !!{structure} कोणत्याही
!!{denumerable} संचावर, त्यात~$\Nat$ वरही, उभारता येते. म्हणून प्रांत
$\Nat$ असलेली अमानक प्रतिमानेही आहेत. अशा~$\Struct{M}$ प्रतिमानांत अर्थातच
किमान काही संख्या आपली नेहमीची भूमिका बजावू शकत नाहीत, कारण असा काही~$k$
असतो की प्रत्येक~$n \in \Nat$ साठी तो~$\Value{\num{n}}{M}$ पेक्षा
वेगळा असतो.
\end{explain}

\begin{defn}
$\Lang{L_A}$ साठीची !!^a{structure}~$\Struct{M}$ ही \emph{संगणनक्षम}
तेव्हाच, जेव्हा $\Domain{M} = \Nat$ असते, $\Assign{\prime}{M}$,
$\Assign{+}{M}$ आणि $\Assign{\times}{M}$ ही संगणनक्षम फलने असतात, आणि
$\Assign{<}{M}$ हा निर्णेय संबंध असतो.
\end{defn}

\begin{ex}\ollabel{ex:comp-model-q}
\olref[mdq]{ex:model-K-of-Q} मधील रचना~$\Struct{K}$ आठवा. तिचा प्रांत
$\Domain{K} = \Nat
\cup \{a\}$ होता आणि अर्थनिर्धारणे पुढीलप्रमाणे होती:
\begin{align*}
  \Assign{\Obj{0}}{K} & = 0\\
  \Assign{\prime}{K}(x) & =
  \begin{cases}
    x+1 & \text{जर $x\in \Nat$}\\
    a & \text{जर $x = a$}
  \end{cases}\\
  \Assign{+}{K}(x, y) & =
  \begin{cases}
    x+y & \text{जर $x$, $y \in\Nat$}\\
    a & \text{इतर प्रसंगी}
  \end{cases}\\
  \Assign{\times}{K}(x, y) & =
  \begin{cases}
    xy & \text{जर $x$, $y \in\Nat$}\\
    0 & \text{जर $x=0$ किंवा $y=0$}\\
    a & \text{इतर प्रसंगी}\\
  \end{cases}\\
  \Assign{<}{K} & =
  \Setabs{\tuple{x,y}}{x, y \in \Nat \text{ आणि } x<y} \cup
  \Setabs{\tuple{x,a}}{x \in \Domain{K}}
\end{align*}
%READERNOTE{OLFOL-040 — गोठवलेल्या इंग्रजी संच-वर्णनात क्रमित जोडीतील x ऐवजी असंबद्ध n वर अट आहे. K च्या घोषित संबंधाशी आणि आधीच्या उदाहरणाशी जुळवून मराठी सूत्रात x बांधला आहे.}
पण $\Domain{K}$ हा !!{denumerable} आहे आणि म्हणून तो~$\Nat$ शी तुल्यबल
आहे. उदाहरणार्थ, $g\colon \Nat \to \Domain{K}$ असे घेऊ, जिथे $g(0) =
a$ आणि $n>0$ साठी $g(n) = n-1$; हे एक !!a{bijection} आहे.
%READERNOTE{OLFOL-039 — गोठवलेल्या इंग्रजी सूत्रात n>0 साठी g(n)=n+1 दिल्यामुळे 0 आणि 1 प्रतिमेत येत नाहीत आणि g द्विएकी राहत नाही. पुढील सर्व स्थानांतरित क्रियांशी सुसंगत असा अपेक्षित नकाशा g(n)=n-1 आहे.}
त्याचा वापर करून नव्या~$\Th{Q}$-प्रतिमान~$\Struct{K'}$ आणि
$\Struct{K}$ यांमध्ये समरूपता मिळवता येते. $\Struct{K'}$ मध्ये
$\Lang{L_A}$ च्या चिन्हांना वेगळी फलने आणि संबंध नेमावे लागतात, कारण
$\Nat$ मधील वेगवेगळे !!{element}s मानक आणि अमानक संख्यांच्या भूमिका
बजावतात.

विशेषतः, आता~$0$ हा लघुतम मानक संख्या नसून~$a$ ची भूमिका बजावतो. आता~$1$
ही लघुतम मानक संख्या आहे. म्हणून $\Assign{\Obj{0}}{K'} = 1$ असे नेमू.
आता उत्तरवर्ती फलनही वेगळे आहे: एखादी मानक संख्या, म्हणजे $n > 0$, दिली
असता ते अजूनही $n+1$ परत करते. पण आता~$0$ हा~$a$ ची भूमिका बजावतो आणि
तो स्वतःचाच उत्तरवर्ती आहे. म्हणून $\Assign{\prime}{K'}(0) = 0$.
बेरीज आणि गुणाकारासाठी याचप्रमाणे
\begin{align*}
\Assign{+}{K'}(x, y) & =
  \begin{cases}
    x+y-1 & \text{जर $x$, $y >0$}\\
    0 & \text{इतर प्रसंगी}
  \end{cases}\\
  \Assign{\times}{K'}(x, y) & =
  \begin{cases}
    1 & \text{जर $x = 1$ किंवा $y = 1$}\\
    xy - x - y + 2 & \text{जर $x$, $y > 1$}\\
    0 & \text{इतर प्रसंगी}\\
  \end{cases}
\end{align*}
आणि $\tuple{x, y} \in \Assign{<}{K'}$ हे तेव्हाच, जेव्हा $x < y$, $x > 0$
आणि $y > 0$ असतात, किंवा $y = 0$ असतो.

ही सर्व फलने नैसर्गिक संख्यांची संगणनक्षम फलने आहेत आणि
$\Assign{<}{K'}$ हा~$\Nat$ वरील निर्णेय संबंध आहे---पण ती~$\Nat$ वरील
उत्तरवर्ती, बेरीज आणि गुणाकार हीच फलने नाहीत, तसेच $\Assign{<}{K'}$ हाही
$\Nat$ वरील~$<$ हाच संबंध नाही.
\end{ex}

\begin{prob}
\olref[mod][mar][mdq]{ex:model-L-of-Q} मधील~$\Struct{L}$ शी समरूपी,
$\Domain{L'} = \Nat$ असलेली !!a{structure}~$\Struct{L'}$ द्या.
\end{prob}

\begin{explain}
\Olref{ex:comp-model-q} यावरून~$\Th{Q}$ ला प्रांत~$\Nat$ असलेली
संगणनक्षम अमानक प्रतिमाने आहेत हे दिसते. पण पुढील निष्कर्ष दाखवतो की
$\Th{PA}$ च्या प्रतिमानांसाठी (आणि म्हणून~$\Th{TA}$ च्या प्रतिमानांसाठीही)
हे खरे नाही.
\end{explain}

\begin{thm}[टेनेनबाउमचे प्रमेय]
$\Struct{N}$ हे~$\Th{PA}$ चे एकमेव संगणनक्षम प्रतिमान आहे.
\end{thm}
  
\end{document}
